વાક્યોના ગણથી એક અર્થમાં એ સંરચનાઓનું વર્ણન થાય છે જેમાં બધાં વાક્યો એકસાથે સાચાં છે. આવી સંરચનાઓ તે ગણના \emph{નિદર્શો} છે. ઊલટું, એક સંરચના અથવા સંરચનાઓનો ગણ આપ્યો હોય તો તે બધા જે વાક્યોના નિદર્શ છે તે વાક્યોનો ગણ જોઈ શકાય છે. તેને સંરચનાનો અથવા સંરચનાઓના ગણનો \emph{સિદ્ધાંત} કહે છે. આવા દરેક ગણમાંથી ફલિત થતું દરેક વાક્ય પહેલેથી તેમાં હોય છે, એટલે તે \emph{બંધ} છે. સામાન્ય રીતે $\Gamma$ ફલિતતા હેઠળ બંધ હોય તો તેને સિદ્ધાંત કહીએ છીએ. $\Gamma$માં $\Delta$માંથી ફલિત થતાં બધાં અને માત્ર એ જ વાક્યો હોય ત્યારે $\Gamma$ એ $\Delta$ વડે \emph{સ્વયંસિદ્ધીકૃત} છે.

ગણિતમાં રેખીય ક્રમોના, જૂથોના અને અંકગણિતના સિદ્ધાંતો મળે છે. પેઆનોની સ્વયંસિદ્ધિઓથી સ્વયંસિદ્ધીકૃત સિદ્ધાંત તેનું એક ઉદાહરણ છે. બીજા વિષયોમાં પણ અગત્યના સિદ્ધાંતો છે. દાખલા તરીકે, સંબંધાત્મક ડેટાબેઝને સિદ્ધાંતો તરીકે જોઈ શકાય છે, અને તત્ત્વમીમાંસામાં ભાગ હોવાના સંબંધ વિશે સ્વયંસિદ્ધીકૃત કરી શકાય એવા સિદ્ધાંતોનો અભ્યાસ થાય છે.

સિદ્ધાંત રચતી વખતે તેની ભાષા ઇચ્છિત બધી વાતો કહી શકે એટલી અભિવ્યક્તિક્ષમ છે કે કેમ અને તે ઇચ્છિત પરિણામો સાબિત કરી શકે એટલો શક્તિશાળી છે કે કેમ તે વિચારવું પડે છે. સંબંધ \emph{વ્યક્ત કરવા} માટે જરૂરી સંખ્યામાં મુક્ત ચલોવાળું સૂત્ર જોઈએ. \emph{ગણસિદ્ધાંત}માં સંબંધ સંકેત તરીકે ફક્ત $\in$ છે. છતાં $x\subseteq y$ને $\lforall[u][(u\in x\lif u\in y)]$ વડે વ્યક્ત કરી શકાય છે. \emph{ઝર્મેલો–ફ્રેન્કેલ ગણસિદ્ધાંત} $\Log{ZFC}$માં થોડાં સ્વયંસિદ્ધિઓ અને $\subseteq$ જેવી વધુ ને વધુ જટિલ વ્યાખ્યાઓની સાંકળ વડે લગભગ દરેક ગાણિતિક વિધાન વ્યક્ત કરી શકાય છે અને લગભગ દરેક ગાણિતિક પ્રમેય સાબિત કરી શકાય છે.
